\documentclass{handout}
\usepackage{handoutSetup}\usepackage{handoutShortcuts}  
\begin{document}

\begin{verbatimwrite}{\jobname.title}
Age Consumption/Income Profiles In Data and Theory 
\end{verbatimwrite}

\handoutNameMake


This handout reproduces figures from \cite{carroll:bslcpih} showing
the profiles of income and consumption by age across different
occupation groups from the 1960-61 {\it Consumer Expenditure Survey}.
The unconstrained perfect foresight/certainty equivalent life cycle
model implies that there should be no relationship between the pattern
of the consumption profile and the income profile, since the shape of
the consumption profile is determined by tastes independently of the
shape of the income profile.

In practice, there seems to be a strong connection between the income
and consumption profiles.  The second set of figures shows the implications
of a ``buffer stock'' model of optimizing life cycle behavior in which 
consumers are modestly impatient, have a standard CRRA utility function
with a precautionary saving motive, and face uncertain income.  
Consumers in this model optimally engage in buffer stock saving 
behavior until 10-15 years before retirement, as a result of their
combination of impatience and prudence.  This model is thus capable
of matching the broad pattern of the data.

\vspace{.2in}



\begin{figure}
\begin{center}
\includegraphics[width=6in]{../Figures/AgeProfilesData}
\end{center}
\end{figure}

\vspace{.5in}

\begin{figure}
\begin{center}
\includegraphics[width=6in]{../Figures/AgeProfilesTheory}
\end{center}
\end{figure}

 
\input handoutBibMake


\end{document}
